\chapter{DishBrain}
\chaptermark{DishBrain}
\label{sec:dishbrain}

\section{Сеть в чашке играет в теннис}

DishBrain --- так авторы назвали стенд, на котором культура живых нейронов, выращенная на чипе с электродами, играет в упрощённый теннис с одной ракеткой~\cite{Kagan2022}. Это самый обсуждаемый опыт с машиной из живых нейронов (обзор таких машин --- в главе~\ref{sec:livingmachine}), и для нашей книги он особенно удобен. Здесь маленькая сеть доступна целиком: каждый вход задаёт экспериментатор, каждый выход записывается. Все вопросы предыдущих глав --- как записан сигнал, кто учит, что видит щуп --- можно задать сети, у которой нет ни глаз, ни мышц, ни \nt{ДОФА}-нейронов. Разберём опыт как инженер разбирает чужой стенд: схема, вход, выход, учитель, результат, споры.

\section{Стенд}

Нейроны берут из коры зародыша мыши --- около $800\,000$ клеток на чип --- или выращивают из стволовых клеток человека, около миллиона на чип. Несколько недель они растут на чипе, прорастают друг в друга и начинают сами выдавать импульсы и пачки. На площадке в $8$~мм$^2$ под ними лежит $26\,000$ платиновых электродов; одновременно можно записывать до $1024$ из них, а стимулировать на практике --- через восемь.

Вокруг культуры замкнута петля (рис.~\ref{fig:dishloop}). Компьютер моделирует игру: мяч летит, отражается от стен, ракетка движется вверх и вниз. Положение мяча превращается в ток на восьми «чувствительных» электродах. Импульсы из двух «двигательных» областей считаются окнами по $10\ms$ и двигают ракетку. После каждого удара или промаха культура получает ещё один стимул --- обратную связь. Окна по $10\ms$ --- это такт компьютера стенда, а не культуры: сама сеть по-прежнему работает асинхронно, как и все сети этой книги.

\begin{figure}[t]
\centering
\begin{tikzpicture}[>=Stealth,
  blk/.style={draw,very thick,rounded corners=3pt,align=center,font=\small,minimum height=1.2cm},
  lab/.style={font=\scriptsize,align=center}]
\node[blk,fill=blue!6,text width=3.4cm] (C) at (0,0) {культура на чипе\\\scriptsize $\sim10^6$ нейронов, $26\,000$ электродов};
\node[blk,fill=orange!10,text width=3.0cm] (G) at (7.2,0) {компьютер:\\игра в теннис};
\draw[->,very thick,blue!60!black] (G.north west) to[bend right=25] node[lab,above] {мяч: \emph{место} --- какой из 8 электродов,\\\emph{частота} --- 4--40 имп/с} (C.north east);
\draw[->,very thick,red!70!black] (C.south east) to[bend right=25] node[lab,below] {ракетка: счёт импульсов\\областей «вверх» и «вниз» за $10\ms$} (G.south west);
\draw[->,thick,densely dashed,green!45!black] (G.east) -- ++(0.9,0) |- ++(0,-2.6) -| node[lab,pos=0.25,below] {отбил: предсказуемый стимул; промах: непредсказуемый} (C.south);
\end{tikzpicture}
\caption{Петля DishBrain. Синяя стрелка --- вход: положение мяча записано кодом местом и частотным кодом (глава~\ref{sec:code}). Красная --- выход: разность счётов двух областей двигает ракетку. Зелёная штриховая --- обратная связь после каждого розыгрыша, единственный «учитель» стенда.}
\label{fig:dishloop}
\end{figure}

\section{Вход и выход}

\emph{Вход} записан сразу двумя кодами главы~\ref{sec:code}. Какой из восьми электродов стимулируется, говорит, на какой высоте мяч: это код местом, как у датчиков главы~\ref{sec:sensors}. Как часто идут импульсы стимуляции, от $4$ до $40$ в секунду, говорит, насколько мяч далеко: это частотный код. Культура не «знает» ни одного из этих кодов заранее: код задан экспериментатором, и сеть должна научиться его читать.

\emph{Выход} читается так же, как в интерфейсах главы~\ref{sec:debugging}. Импульсы одной двигательной области толкают ракетку вверх, другой --- вниз; в простейшей записи за каждое окно
\begin{empheq}[box=\keybox]{equation}
\Delta p \;=\; c\,\bigl(n_\uparrow - n_\downarrow\bigr),
\label{eq:dishreadout}
\end{empheq}
где $n_\uparrow$, $n_\downarrow$ --- числа импульсов двух областей за $10\ms$, а $c$ --- коэффициент стенда. Это двухпроводный код главы~\ref{sec:glutcomputer}: знак движения несёт не знак сигнала, а то, какой из двух проводов громче.

Посмотрим на шум такого выхода. Если область выдаёт в сумме $R$ импульсов в секунду, за окно $T=10\ms$ их в среднем $RT$, и по~\eqref{eq:rateerror} относительный разброс равен $1/\sqrt{RT}$. При $R=1000$ импульсов в секунду это около $30$ процентов в каждом окне; при $R=100$ --- больше ста. Ракетка неизбежно дрожит, и обучиться сеть может лишь одним способом --- сдвинуть \emph{среднюю} разность счётов в нужную сторону. Это те же законы, что ограничивают любой интерфейс главы~\ref{sec:debugging}, только теперь по обе стороны петли.

\section{Учитель без дофамина}

Самое интересное в стенде --- учитель. В культуре нейронов коры нет \nt{ДОФА}-нейронов, и никакого сигнала «лучше, чем ожидалось» стенд не посылает. Обратная связь другая. Отбила сеть мяч --- на все восемь чувствительных электродов сразу подаётся короткий \emph{предсказуемый} стимул: $100$ импульсов в секунду в течение $0.1\,\text{с}$. Промахнулась --- четыре секунды идёт стимул, который авторы называют \emph{непредсказуемым}: импульсы по $150$~мВ, $5$ в секунду~\cite{Kagan2022}.

Авторы исходили из гипотезы, что сеть действует так, чтобы её вход становился предсказуемым~\cite{Friston2010}. Для инженера смысл простой: у культуры есть ровно один способ избавиться от четырёх секунд случайного шума --- отбить мяч. Та же идея встречалась нам в экономии импульсов главы~\ref{sec:code} и в предсказании кадра в главе~\ref{sec:recognition}.

Но нужна ли здесь именно предсказуемость? Возможен и более грубый механизм. Пусть непредсказуемый стимул после промаха просто \emph{перемешивает} недавние изменения в сети, а спокойный после удара их не трогает. Опишем настройку сети одним вектором $\boldsymbol\psi$, и пусть в настройке $\boldsymbol\psi$ сеть промахивается с вероятностью $P_{\text{пр}}(\boldsymbol\psi)$. Если толчки симметричны --- из $\boldsymbol\psi$ в $\boldsymbol\psi'$ так же вероятно, как обратно, --- то в установившемся режиме поток уходов из каждой настройки уравновешен потоком приходов, и доля времени, которую сеть проводит в настройке $\boldsymbol\psi$, равна
\begin{empheq}[box=\keybox]{equation}
\rho(\boldsymbol\psi) \;\propto\; \frac{1}{P_{\text{пр}}(\boldsymbol\psi)} .
\label{eq:dishdensity}
\end{empheq}
Настройка, которая промахивается вдесятеро реже, удерживается вдесятеро дольше. Никто не сообщает сети, в какую сторону менять веса: хорошие настройки просто реже разрушаются.

Мы проверили это на модели, сведя культуру к двум числам: ракетка приходит в точку $a\,y+b$, когда мяч приходит на высоту $y$, и отбивает, если промахнулась меньше чем на $0.1$. После промаха оба числа получают случайный толчок, после удара остаются прежними. Доля отбитых мячей за $2000$ розыгрышей выросла с $0.21$ в первых двухстах до $0.38$ в последних пятистах (сорок «культур», рис.~\ref{fig:dishmodel}). Если толчок давать после \emph{каждого} розыгрыша, сведений в обратной связи нет, и доля остаётся около $0.12$; если толчков нет вовсе, она стоит на $0.19$. Опыт на стенде устроен так же: улучшение было только при полной обратной связи, а при тишине вместо шума после промаха и без обратной связи его не было~\cite{Kagan2022}.

\begin{figure}[t]
\centering
\begin{tikzpicture}[x=1cm,y=5cm]
\draw[->,gray!70!black] (0,0) -- (10.6,0);
\draw[->,gray!70!black] (0,0) -- (0,0.5) node[above,font=\small,black] {доля отбитых};
\foreach \v in {0.1,0.2,0.3,0.4}{\draw[gray!60!black] (-0.06,\v) -- (0.06,\v); \node[left,font=\scriptsize] at (0,\v) {\v};}
\foreach \x/\f/\l/\lab in {1.8/0.21/0.38/{шум после\\промаха},5.2/0.13/0.11/{шум после\\каждого розыгрыша},8.6/0.19/0.19/{без шума}}{
  \fill[blue!35] (\x-0.8,0) rectangle (\x-0.05,\f);
  \fill[red!60!black] (\x+0.05,0) rectangle (\x+0.8,\l);
  \node[below,font=\scriptsize,align=center] at (\x,-0.01) {\lab};
  \node[above,font=\scriptsize] at (\x-0.425,\f) {\f};
  \node[above,font=\scriptsize] at (\x+0.425,\l) {\l};}
\fill[blue!35] (7.4,0.47) rectangle (7.7,0.495); \node[right,font=\scriptsize] at (7.75,0.482) {первые 200};
\fill[red!60!black] (7.4,0.41) rectangle (7.7,0.435); \node[right,font=\scriptsize] at (7.75,0.422) {последние 500};
\end{tikzpicture}
\caption{Модель «перемешивание при промахе» (\texttt{dishbrain\_model.py}): доля отбитых мячей в начале и в конце $2000$ розыгрышей, среднее по сорока моделям-культурам. Обучение идёт, только если толчок следует за промахом, но не за ударом, --- как и в опыте, где сеть училась лишь при полной обратной связи.}
\label{fig:dishmodel}
\end{figure}

\begin{keyblock}
\begin{proposition}[Учитель-перемешиватель]
\label{prop:dishscramble}
Если неудача встряхивает сеть, а удача оставляет её в покое, сеть сама дрейфует к настройкам, которые реже ошибаются: время, проведённое в настройке, обратно пропорционально вероятности неудачи~\eqref{eq:dishdensity}. Для такого обучения не нужны ни сигнал награды, ни его знак, ни предсказуемость как таковая --- достаточно, чтобы обратная связь после удачи и после неудачи различалась. Изменение поведения культуры измерено; какой механизм работает в чашке --- предсказание входа или перемешивание, --- не установлено.
\end{proposition}
\end{keyblock}

Сравните с главой~\ref{sec:dofa}. Там шум тоже служил поиском, но у дофамина был знак: ошибка $\delta$ говорила, в какую сторону сдвинуть веса, и шаг шёл по градиенту~\eqref{eq:perturbation}. Здесь знака нет --- есть только «оставить» или «встряхнуть». Такой поиск грубее и медленнее, но требует от сети меньше: ни \nt{ДОФА}-нейронов, ни критика, ни следа, растянутого на секунды.

\section{Что измерили и о чём спорят}

Каждая игра длилась $20$ минут; сравнивали первые $5$ минут с последними $15$. У культур с полной обратной связью серии отбитых мячей стали длиннее, а мгновенных промахов --- меньше; в контролях без клеток, без игры и с компьютерной «ракеткой» ничего подобного не было. Культуры из клеток человека в среднем отбивали дольше мышиных. Улучшение статистически надёжно --- на девяти мышиных и одиннадцати человеческих культурах, в сотне с лишним игр каждой группы, --- но невелико: сеть не стала хорошим игроком, она стала немного лучше случайного~\cite{Kagan2022}. В следующей работе той же группы нашли, что во время игры активность культуры приближается к \emph{критическому} режиму --- границе между затуханием и лавиной, той самой жизни на грани, что в главе~\ref{sec:gaba}, --- и чем ближе, тем лучше игра~\cite{Habibollahi2023}.

Главный спор вызвали не цифры, а слова. Заголовок статьи утверждал, что нейроны в чашке «проявляют разумность» (sentience), и группа исследователей возразила: изменение поведения в замкнутой петле --- ещё не разум и не ощущение, а выводы нужно формулировать скромнее~\cite{Balci2023}. С инженерной точки зрения открыты два вопроса, и оба --- о механизме. Первый: учится ли сеть, то есть меняются ли её веса надолго, или обратная связь лишь сдвигает её текущее состояние? Второй: нужна ли именно предсказуемость, или хватает любого различия между откликом на удачу и на неудачу, как в утверждении~\ref{prop:dishscramble}? Стенд, на котором доступен каждый электрод, позволяет ответить на оба --- и это, пожалуй, его главная ценность.

\section{Спецификация стенда: как его воспроизвести}
\label{sec:dishspec}

Ниже собрано всё, что нужно, чтобы собрать такой стенд заново: оборудование, петля, кодирование входа, чтение выхода, обратная связь и протокол опыта. Каждый параметр помечен источником. «Статья» --- значение взято из текста работы~\cite{Kagan2022}. «Выбор» --- в статье его нет, и для воспроизведения его придётся задать самому; мы предлагаем значение по умолчанию и говорим, как его проверить.

\subsection*{Оборудование и культура}

\begin{center}
\footnotesize
\setlength{\tabcolsep}{4pt}
\renewcommand{\arraystretch}{1.15}
\begin{tabular}{>{\raggedright}p{4.0cm}>{\raggedright}p{6.9cm}>{\raggedright\arraybackslash}p{1.6cm}}
\toprule
Параметр & Значение & Источник \\
\midrule
чип & матрица MaxOne: $26\,000$ платиновых электродов на площадке $8$~мм$^2$ & статья \\
запись & до $1024$ каналов одновременно, $20\,000$ отсчётов в секунду & статья \\
стимуляция & до $32$ электродов технически, в опыте --- $8$ независимых & статья \\
клетки & кора зародыша мыши; нейроны человека из стволовых клеток (два способа выращивания); клетки HEK293T (не нейроны) для контроля & статья \\
посев & около $10^6$ клеток на чип & статья \\
возраст культуры к опыту & мышь: с $\sim10$ суток; человек: $14$--$24$ суток или $\sim80$ суток в зависимости от способа выращивания & статья \\
\bottomrule
\end{tabular}
\end{center}

\subsection*{Петля и время}

Петля работает окнами по $10\ms$ ($200$ отсчётов): за окно считаются импульсы на электродах двигательных областей, игра обновляется, и выдаётся очередная стимуляция. Задержка от импульса в культуре до ответного стимула --- около $5\ms$. Во время игры сигнал каждого электрода проходит фильтр верхних частот Бесселя 2-го порядка с частотой среза $100$~Гц; модуль сигнала сглаживается фильтром нижних частот Бесселя 1-го порядка с частотой $1$~Гц, и порог импульса пропорционален этому сглаженному модулю. Порог в шесть стандартных отклонений фона статья указывает для отдельных измерений активности, а не для игры. Чтобы стимуляция не засчитывалась как ответ культуры, все сигналы гасятся, когда одновременно приходит больше $15$ крупных, выше $75$~мВ, импульсов. Всё это --- статья.

\subsection*{Кодирование мяча}

\begin{center}
\footnotesize
\setlength{\tabcolsep}{4pt}
\renewcommand{\arraystretch}{1.15}
\begin{tabular}{>{\raggedright}p{3.4cm}>{\raggedright}p{7.5cm}>{\raggedright\arraybackslash}p{1.6cm}}
\toprule
Параметр & Значение & Источник \\
\midrule
чувствительная область & $8$ электродов, расположенных так, чтобы соседние электроды означали соседние положения мяча & статья \\
код местом & номер электрода $k=1,\dots,8$ задаёт положение мяча относительно центра ракетки & статья \\
какая координата & вертикальное положение мяча; для воспроизведения --- относительно ракетки, $k=\lceil 8(y-p+\tfrac12)\rceil$ при $y-p\in[-\tfrac12,\tfrac12]$ & статья; формула --- выбор \\
код частотой & частота стимуляции от $4$ до $40$~имп/с несёт расстояние до ракетки & статья \\
направление шкалы & $4$~имп/с, когда мяч у противоположной стены, и линейно до $40$~имп/с, когда он у стены ракетки: $f=4+36\,(1-x/L)$, $x$ --- расстояние до стены ракетки, $L$ --- ширина корта & статья \\
форма импульса & короткий прямоугольный двухфазный: сначала положительная, потом отрицательная фаза & статья \\
амплитуда & $75$~мВ; выбрана так, чтобы не заставлять срабатывать заторможенные клетки & статья \\
длительность фаз & не указана; взять стандартную настройку системы стимуляции и не менять её в течение опыта & выбор \\
\bottomrule
\end{tabular}
\end{center}

\subsection*{Чтение ракетки}

В статье две двигательные области: импульсы в первой двигают ракетку вверх, во второй --- вниз; вклад областей взвешивается, чтобы учесть разную плотность клеток и разную активность~\cite{Kagan2022}. Точная формула не приведена. Для воспроизведения берём правило~\eqref{eq:dishreadout}, $\Delta p=c\,(n_\uparrow-n_\downarrow)$ за каждое окно $10\ms$, с весом, уравнивающим области по их средней активности в покое (выбор): $n_\uparrow$ и $n_\downarrow$ делятся на средний счёт своей области за окно, измеренный перед игрой. Коэффициент $c$ выбираем так, чтобы разница в один средний счёт сдвигала ракетку на $1\%$ высоты корта за окно (выбор); тогда постоянный перевес одной области проводит ракетку через весь корт за $1$~с.

Расположение областей на чипе в тексте статьи координатами не задано; есть только схема. Для воспроизведения достаточно трёх непересекающихся групп электродов: восемь чувствительных в центре и по группе записывающих электродов с двух сторон, одна --- «вверх», другая --- «вниз» (выбор).

\subsection*{Игра}

Размеры корта, длина ракетки и скорость мяча в статье не указаны; сказано только, что мяч летит с постоянной скоростью и отражается от стен. Для воспроизведения (выбор): корт $1\times1$, ракетка на правой стенке длиной $0.2$ (попадание при $|y-p|<0.1$, как в модели~\texttt{models/dishbrain\_model.py}), мяч пересекает корт за $1$~с. Промах без единого отбитого мяча в розыгрыше считается «эйсом»; розыгрыш длиннее трёх ударов подряд --- «длинным» (статья).

\subsection*{Обратная связь}

\begin{center}
\footnotesize
\setlength{\tabcolsep}{4pt}
\renewcommand{\arraystretch}{1.15}
\begin{tabular}{>{\raggedright}p{2.6cm}>{\raggedright}p{8.3cm}>{\raggedright\arraybackslash}p{1.6cm}}
\toprule
Событие & Что подаётся & Источник \\
\midrule
удар & предсказуемый стимул: все $8$ чувствительных электродов одновременно, $75$~мВ, $100$~имп/с, $100\ms$; обычная стимуляция мяча на это время прерывается & статья \\
промах & непредсказуемый стимул: $150$~мВ, $5$~имп/с, $4$~с, в случайные моменты на случайно выбранные из $8$ электродов; затем $4$~с без стимуляции, и мяч стартует в случайном направлении & статья \\
закон случайности & не указан; для воспроизведения --- моменты импульсов как пуассоновский поток со средней частотой $5$~имп/с & выбор \\
\bottomrule
\end{tabular}
\end{center}

\subsection*{Протокол и контроль}

Один сеанс длится $20$~мин; сравнивают первые $5$ и последние $15$~мин (статья). Три меры результата: средняя длина розыгрыша, доля эйсов и доля длинных розыгрышей. Условия опыта (статья):
\begin{itemize}
\item \emph{стимул} --- полная петля, как описано выше;
\item \emph{тишина} --- вместо стимула удара или промаха столько же времени без всякой стимуляции, затем мяч стартует в случайном направлении;
\item \emph{без обратной связи} --- мяч после промаха просто отражается и летит дальше, стимуляция положения мяча продолжается;
\item \emph{покой} --- игра идёт и ракетка движется по активности, но стимуляции нет совсем;
\item \emph{контроли} --- чип только со средой; клетки HEK293T; «культура в компьютере», симуляция вместо клеток.
\end{itemize}
Объём выборки в основной серии: $9$ культур мыши ($101$ сеанс), $11$ культур человека ($138$), $6$ чипов со средой ($80$), $20$ культур в покое ($42$), $3$ симуляции ($38$), всего $399$ сеансов. В серии про обратную связь --- $486$ сеансов за $3$ дня по $3$ сеанса в день, условия «стимул», «тишина» и «без обратной связи» ($n=27$, $17$ и $15$). Рост средней длины розыгрыша от первых $5$ минут к последним $15$ получен только у живых культур и только в условии «стимул»~\cite{Kagan2022}.

\subsection*{Цикл стенда по шагам}

Каждые $10\ms$ компьютер стенда делает следующее.
\begin{enumerate}
\item Считывает число импульсов $n_\uparrow$, $n_\downarrow$ в двух двигательных областях за прошедшее окно и нормирует их на средний счёт области в покое.
\item Сдвигает ракетку на $\Delta p=c\,(n_\uparrow-n_\downarrow)$ и ограничивает её краями корта.
\item Сдвигает мяч на один шаг; отражает его от верхней, нижней и левой стенок.
\item Если мяч дошёл до правой стенки: при $|y-p|<0.1$ --- удар, счёт розыгрыша растёт на один, подаётся стимул удара ($100\ms$); иначе --- промах, розыгрыш записывается, подаётся стимул промаха ($4$~с), затем $4$~с паузы, и мяч стартует заново.
\item Иначе выбирает электрод $k$ по вертикальному положению мяча и частоту $f$ по расстоянию до стены ракетки и подаёт на электрод $k$ двухфазные импульсы $75$~мВ с частотой $f$ до следующего окна.
\end{enumerate}
Этого достаточно, чтобы собрать стенд, совместимый с опубликованным, и проверить главный вопрос главы: учит ли культуру предсказуемость или простая разница между ответом на удар и на промах. Для этого к трём условиям статьи стоит добавить четвёртое, которого в ней нет: после промаха --- \emph{предсказуемый} стимул той же длительности и энергии, что и непредсказуемый. Если культура учится и так, дело в разнице, а не в предсказуемости.


\section{Итог}

\begin{table}[t]
\centering
\footnotesize
\setlength{\tabcolsep}{4pt}
\renewcommand{\arraystretch}{1.15}
\begin{tabular}{>{\raggedright}p{2.1cm}>{\raggedright}p{4.2cm}>{\raggedright}p{1.9cm}>{\raggedright\arraybackslash}p{4.6cm}}
\toprule
Узел стенда & Как устроен & Глава книги & Что известно \\
\midrule
вход & место (какой из 8 электродов) и частота (4--40 имп/с) & \ref{sec:code}, \ref{sec:sensors} & задано экспериментатором \\
выход & разность счётов двух областей за $10\ms$~\eqref{eq:dishreadout} & \ref{sec:glutcomputer}, \ref{sec:debugging} & задано экспериментатором; шум по~\eqref{eq:rateerror} \\
учитель & предсказуемый стимул после удара, непредсказуемый после промаха; дофамина нет & \ref{sec:dofa} & улучшение при полной обратной связи измерено; без неё его нет \\
механизм & предсказание входа или перемешивание при промахе~\eqref{eq:dishdensity} & \ref{sec:dofa}, \ref{sec:recognition} & не установлен; оба объясняют результат \\
режим сети & приближение к критическому режиму во время игры & \ref{sec:gaba} & связь с качеством игры измерена \\
«разумность» & утверждение авторов & --- & не показана; оспорена \\
\bottomrule
\end{tabular}
\caption{DishBrain глазами инженера.}
\label{tab:dishbrain}
\end{table}

DishBrain --- машина из живых нейронов в самом простом виде: вход задан кодом, выход прочитан счётом, учитель --- различие между откликом на удачу и на неудачу (таблица~\ref{tab:dishbrain}). Сеть без глаз, мышц и дофамина чуть-чуть научилась играть, и уже одно это превращает её в испытательный стенд для идей книги. Каким бы ни оказался механизм --- предсказание, перемешивание или что-то третье, --- ответ будет найден так, как советует глава~\ref{sec:debugging}: щупом, тестовым сигналом и выключением частей на машине, которую наконец видно целиком.

\section{Задачи}

\begin{exercise}[\lvlA]
\label{ex:22:1}
\emph{Сколько нужно сдвинуть счёт.} Две двигательные области выдают вместе $R_\uparrow+R_\downarrow=2000$ импульсов в секунду, потоки пуассоновские и независимые.
(а)~Каков относительный разброс счёта одной области за окно $10\ms$ при $R=1000$ и $R=100$ импульсов в секунду?
(б)~Мяч летит до ракетки $1$~с, и за это время смещения~\eqref{eq:dishreadout} складываются. Какой наименьшей разности $R_\uparrow-R_\downarrow$ достаточно, чтобы среднее смещение ракетки вдвое превышало его разброс? Какую долю от суммарной частоты это составляет?
\end{exercise}

\begin{exercise}[\lvlA]
\label{ex:22:2}
\emph{Сколько розыгрышей нужно, чтобы увидеть обучение.} Считайте розыгрыши независимыми, а долю отбитых мячей --- биномиальной.
(а)~Сколько розыгрышей нужно в каждом из двух периодов, чтобы рост доли отбитых с $0.20$ до $0.25$ превышал два стандартных отклонения разности?
(б)~То же для роста с $0.21$ до $0.38$, как в модели рис.~\ref{fig:dishmodel}.
(в)~Почему в опыте розыгрыши одной культуры нельзя считать независимыми и зачем поэтому нужны десятки культур?
\end{exercise}

\begin{exercise}[\lvlB]
\label{ex:22:3}
\emph{Вывод~\eqref{eq:dishdensity}.} Пусть настройка $\boldsymbol\psi$ принимает значения из конечного множества. В каждом розыгрыше сеть промахивается с вероятностью $P(\boldsymbol\psi)>0$ и тогда переходит в $\boldsymbol\psi'$ с вероятностью $K(\boldsymbol\psi,\boldsymbol\psi')=K(\boldsymbol\psi',\boldsymbol\psi)$; при ударе остаётся на месте.
(а)~Докажите, что $\rho(\boldsymbol\psi)\propto1/P(\boldsymbol\psi)$ --- стационарное распределение, и укажите условие его единственности.
(б)~Докажите, что установившаяся доля промахов равна среднему гармоническому $P$ по всем настройкам, и что она не больше доли промахов при отсутствии обратной связи. Когда достигается равенство?
(в)~Пример: две настройки с $P=0.9$ и $P=0.1$. Найдите обе доли промахов.
(г)~Что произойдёт, если $P=0$ в некоторых настройках?
\end{exercise}

\begin{exercise}[\lvlB]
\label{ex:22:4}
\emph{Поглощающая область модели.} В модели рис.~\ref{fig:dishmodel} ракетка приходит в $p=\min\bigl(1,\max(0,ay+b)\bigr)$, мяч --- на высоту $y\sim U(0,1)$, удар --- если $|p-y|<h$, $h=0.1$.
(а)~Покажите, что ограничение $p$ отрезком $[0,1]$ никогда не увеличивает промах.
(б)~Докажите, что при $|b|<h$ и $|a-1+b|<h$ сеть отбивает все мячи, и найдите площадь этой области на плоскости $(a,b)$. Почему такие настройки --- поглощающие?
(в)~Найдите численно среднюю долю отбитых мячей при начальном распределении модели ($a\sim U(-1,1)$, $b\sim U(0,1)$) и сравните со столбцом «без шума» рис.~\ref{fig:dishmodel}.
\end{exercise}

\begin{exercise}[\lvlB]
\label{ex:22:5}
\emph{Моделирование с \texttt{dishbrain\_model.py}.} Работайте в копии модели; усредняйте по $200$ культурам (затравки $0$--$199$).
(а)~Увеличьте число розыгрышей до $20\,000$. Какой будет доля отбитых в последних $500$? Какая доля культур к концу оказалась в поглощающей области задачи~\ref{ex:22:4}? Почему доля отбитых не стремится к единице?
(б)~Ограничьте настройку прямоугольником $a\in[-1,2]$, $b\in[-1,1]$ (обрезайте после каждого толчка) и повторите для $2000$ и $20\,000$ розыгрышей. Объясните разницу с помощью задачи~\ref{ex:22:3}.
(в)~Для $2000$ розыгрышей постройте долю отбитых в последних $500$ как функцию силы толчка \texttt{sigma} от $0.02$ до $0.64$. Где оптимум и почему он порядка $h$?
\end{exercise}

\begin{exercise}[\lvlB]
\label{ex:22:6}
\emph{Проектирование входного кода.} Ракетка отбивает мяч, если ошибка по высоте меньше $h=0.1$. Сеть успевает прочитать вход за $T=0.25$~с; частота стимуляции ограничена $4$--$40$ импульсами в секунду; по главе~\ref{sec:code} до частоты $r$ за время $T$ различимо около $\sqrt{rT}$ уровней.
(а)~Сколько уровней высоты нужно различать? Хватает ли кода местом на восьми электродах?
(б)~Можно ли было бы передать высоту частотой на одном электроде? Какая предельная частота для этого понадобилась бы?
(в)~Сколько бит о расстоянии до мяча несёт частотный код стенда? Предложите распределение двух величин между кодами при том же числе импульсов стимуляции и обоснуйте его.
\end{exercise}

\begin{exercise}[\lvlC]
\label{ex:22:7}
\emph{Исследование: предсказание или перемешивание?} Глава оставляет открытым, что работает в чашке.
(а)~Обобщите модель на $d$ настраиваемых чисел (например, ракетка $p=\sum_{i=0}^{d-1}a_iy^i$) и численно найдите, как число розыгрышей до доли отбитых $0.5$ растёт с $d$ при перемешивании. Сравните с правилом, в котором знак ошибки известен (глава~\ref{sec:dofa}, \eqref{eq:perturbation}). Какая зависимость от $d$ ожидается в каждом случае?
(б)~Предложите опыт на стенде, в котором две гипотезы дают противоположные предсказания. Например: после промаха подавать \emph{предсказуемый}, но сильный стимул, а после удара --- непредсказуемый, но слабый. Что предскажет каждая гипотеза?
(в)~Оцените по задаче~\ref{ex:22:2}, сколько культур и розыгрышей нужно, чтобы опыт (б) дал надёжный ответ при ожидаемой разнице долей около $0.05$.
\end{exercise}
